% Budget : 17 pages — voir rapport/PLAN.md
%
% Figures : exports PDF de draw.io, laisses dans diagrams/drawio/ et atteints par le
% \graphicspath du preambule — uc-global.pdf et class-domain.pdf.
% Toujours exporter en PDF plutot qu'en PNG : vectoriel, net a l'impression, bien plus leger.
\chapter{Sprint 0 — Initialisation, besoins et étude technique}
\label{ch:sprint0}

\section{Introduction}

Le chapitre précédent a posé le cadre : une entreprise de services numériques, un besoin de suivi de
livraison adossé à des progiciels de gestion hétérogènes, et une méthode de travail itérative. Reste
à transformer ce cadre en un plan de travail exécutable.

C'est l'objet du sprint~0. Contrairement aux cinq sprints qui le suivent, il ne livre aucune
fonctionnalité~: il produit les décisions sans lesquelles les autres ne pourraient pas commencer.
Qui utilise le système et avec quels droits, ce qu'il doit savoir faire, dans quel ordre le
construire, et avec quels outils.

Ce chapitre restitue ces quatre travaux dans cet ordre. Il se termine sur trois exigences qui, sans
être des fonctionnalités, ont pesé plus lourd que toutes les autres sur les choix qui ont suivi.

\section{Identification des acteurs}

Un acteur est un rôle joué par une personne ou un système extérieur qui interagit avec la solution.
Le découpage retenu ne décalque pas l'organigramme de l'entreprise cliente~: il suit les
\emph{permissions réellement distinctes}. Deux fonctions qui accèdent aux mêmes écrans et déclenchent
les mêmes opérations ne forment qu'un seul acteur, quel que soit leur intitulé sur la fiche de poste.

Ce principe a une conséquence concrète, détaillée au chapitre~\ref{ch:socle}~: le système distingue
trois rôles internes et non quatre, parce que superviser et exploiter ne mobilisent pas les mêmes
droits, tandis que deux intitulés qui exploitent de la même façon ne justifient pas deux rôles
séparés.

\begin{table}[H]
\centering
\caption{Acteurs du système et périmètre de chacun}
\label{tab:acteurs}
\begin{tabular}{L{3.1cm}L{5.2cm}L{6.2cm}}
\hline
\rowcolor{isimsblue!10}
\textbf{Acteur} & \textbf{Description} & \textbf{Périmètre d'action} \\
\hline
Administrateur & Responsable du paramétrage de l'installation pour son entreprise. &
Comptes et rôles, correspondance des champs avec le progiciel, référentiel des chauffeurs et des
véhicules. Hérite de toutes les actions du dispatcheur. \\
\hline
Dispatcheur & Il transforme au quotidien les commandes en tournées et arbitre les aléas. & Import des commandes, planification, affectation, validation, traitement des exceptions et
des retours, comptage et arbitrage des remises de caisse, zones et dépôts. \\
\hline
Responsable & Supervise l'activité sans intervenir dessus. &
Consultation des livraisons, des tournées, des rapports, des espèces en circulation et du journal
d'audit. Aucune action de modification. \\
\hline
Livreur & Exécute les tournées sur le terrain depuis l'application mobile. &
Consultation de sa tournée du jour, changements de statut, preuve de livraison, déclaration d'échec,
encaissement, transfert de garde. \\
\hline
Destinataire & Client final du chargeur. &
Consultation de l'avancement de son colis par un lien public, sans compte ni authentification. \\
\hline
\end{tabular}
\end{table}

Deux systèmes extérieurs complètent ce tableau sans être des acteurs au sens strict, puisqu'ils
n'ont ni intention ni objectif propre~: le \textbf{progiciel de gestion} de l'entreprise cliente
— Odoo ou ERPNext —, source des commandes et destinataire des statuts, et le \textbf{serveur
d'identité} Keycloak, qui authentifie les utilisateurs internes. Ils apparaissent dans le diagramme
des cas d'utilisation sous forme de composants stéréotypés, à droite de la frontière du système.

\section{Analyse et spécification des besoins}
\label{sec:besoins}

\subsection{Besoins fonctionnels}

Les besoins fonctionnels décrivent ce que le système doit savoir faire. Ils sont regroupés en huit
domaines, qui serviront ensuite de découpage au \emph{product backlog}.

\paragraph*{Gestion des identités et des accès.} Le système doit authentifier les utilisateurs
internes, leur attribuer un rôle, et n'autoriser chaque opération que si le rôle porte la permission
correspondante. Il doit également héberger plusieurs entreprises clientes sur une même installation,
sans qu'aucune ne puisse atteindre les données d'une autre.

\paragraph*{Intégration au progiciel de gestion.} Le système doit importer les commandes prêtes à
être livrées, puis renvoyer les statuts de livraison, les quantités réellement livrées et les preuves
associées. Cette intégration doit fonctionner avec des progiciels différents, et avec des versions
différentes d'un même progiciel, sans réécriture du cœur applicatif. L'administrateur doit pouvoir
établir lui-même la correspondance entre les champs de son progiciel et ceux du système.

\paragraph*{Référentiel logistique.} Le système doit gérer les chauffeurs, les véhicules et les zones
de couverture. Les dépôts, en revanche, ne sont pas saisis~: ils reflètent les entrepôts du
progiciel et sont récupérés par synchronisation, afin d'éviter un double référentiel qui divergerait
au premier changement.

\paragraph*{Planification et affectation.} Le dispatcheur doit pouvoir construire une tournée à partir
d'un ensemble de livraisons, y affecter un chauffeur et un véhicule, réordonner les arrêts, demander
une optimisation automatique du parcours, puis valider la tournée pour la rendre visible au
chauffeur.

\paragraph*{Exécution terrain.} Le chauffeur doit consulter sa tournée, déclarer son départ, marquer
son arrivée à chaque arrêt, livrer ou déclarer un échec motivé, et collecter une preuve de livraison
comportant le nom du destinataire, une photographie du colis et une photographie du bon signé. Un
transfert de garde entre deux chauffeurs doit être possible en cours de journée, avec confirmation
du receveur.

\paragraph*{Encaissement et remise de caisse.} Lorsque la commande est payable à la livraison, le
chauffeur doit enregistrer le montant encaissé et son mode de règlement. En fin de tournée, il
déclare le total remis~; une autre personne compte la somme reçue et arbitre l'écart éventuel.

\paragraph*{Suivi et traitement des aléas.} Le dispatcheur doit voir l'avancement des livraisons en
temps réel, être alerté des retards et des échecs, réaffecter une livraison bloquée, et enregistrer
les retours de marchandise. Le destinataire doit pouvoir suivre son colis par un lien public.

\paragraph*{Pilotage et exploitation.} Le responsable doit disposer d'indicateurs — taux de réussite,
motifs d'échec, respect des délais, performance par chauffeur et par zone — et pouvoir exporter des
documents. L'administrateur doit disposer d'une console de santé qui expose l'état de l'intégration et
permet de relancer une synchronisation en échec.

\subsection{Besoins non fonctionnels}

Les besoins non fonctionnels ne décrivent pas des fonctions mais des qualités attendues du système.
Trois d'entre eux ont pesé plus lourd que les autres sur l'architecture et sont signalés comme tels.

\begin{table}[H]
\centering
\caption{Besoins non fonctionnels}
\label{tab:besoins-nf}
\begin{tabular}{L{3.4cm}L{11.1cm}}
\hline
\rowcolor{isimsblue!10}
\textbf{Besoin} & \textbf{Exigence} \\
\hline
\textbf{Agnosticité ERP} & Le cœur métier ne doit connaître aucun progiciel en particulier. Ajouter
un progiciel, ou prendre en charge une nouvelle version, ne doit pas modifier le code métier.
\emph{Structurant.} \\
\hline
\textbf{Isolation par client} & Les données de chaque entreprise cliente doivent être séparées de
façon structurelle, et non par un simple filtre applicatif qu'un oubli suffirait à contourner.
\emph{Structurant.} \\
\hline
\textbf{Continuité hors ligne} & L'application chauffeur doit rester utilisable sans réseau. Les
actions réalisées hors couverture doivent être rejouées à la reconnexion sans produire de doublon.
\emph{Structurant.} \\
\hline
Sécurité & Communications chiffrées, jetons d'accès à durée limitée, contrôle des droits appliqué à
la fois à la passerelle et dans chaque service. \\
\hline
Traçabilité & Toute opération sensible doit laisser une trace horodatée et attribuée à son auteur,
consultable a posteriori. \\
\hline
Temps réel & Un changement d'état survenu sur le terrain doit se refléter sur le poste du
dispatcheur sans rechargement manuel. \\
\hline
Performance & Les listes doivent être renvoyées par pages plutôt qu'en bloc, et tout critère de
recherche exposé à l'utilisateur doit être indexé. \\
\hline
Portabilité & L'ensemble doit se déployer sur une machine unique par conteneurs, sans installation
manuelle de dépendances. \\
\hline
Ergonomie & Interface d'administration en français, arabe et anglais~; application chauffeur
conçue pour un usage à une main. \\
\hline
\end{tabular}
\end{table}

\subsection{Diagramme global des cas d'utilisation}

Le diagramme de la figure~\ref{fig:uc-global} synthétise les besoins fonctionnels sous forme de cas
d'utilisation et les rattache à leurs acteurs. Il reste volontairement au premier niveau~: le
système compte une trentaine de cas, les représenter tous rendrait la figure illisible. Le détail
est porté par les diagrammes propres à chaque sprint, aux chapitres~\ref{ch:socle} à~\ref{ch:ops}.

\begin{figure}[H]
\centering
% Diagramme bien plus haut que large : le borner à la seule largeur donnerait 227 mm de hauteur
% et ferait déborder la page. On borne les deux dimensions et keepaspectratio retient la plus
% contraignante.
\includegraphics[width=\textwidth,height=0.86\textheight,keepaspectratio]{uc-global}
\caption{Diagramme global des cas d'utilisation}
\label{fig:uc-global}
\end{figure}

Trois éléments de lecture méritent d'être soulignés.

\textbf{Une seule généralisation relie les acteurs}~: l'administrateur est un dispatcheur qui dispose
en plus du paramétrage, et ses permissions incluent effectivement toutes celles du dispatcheur. La
flèche évite de redessiner une douzaine de traits.

Le responsable, lui, reste rattaché à ses propres cas. Il croise le dispatcheur sur les consultations
sans jamais le contenir ni en dériver~: il est seul à accéder au journal d'audit, et n'a aucune des
permissions d'action. Superviser suppose de pouvoir relire ce que font ceux qui opèrent, sans
pouvoir opérer soi-même~; c'est pourquoi les rôles sont déclarés à plat dans l'annuaire, chacun
portant sa liste explicite, plutôt qu'en chaîne d'héritage. Le chapitre~\ref{ch:socle} détaille ce
modèle.

Les liens \emph{«~include~»} vers le cas \emph{S'authentifier} ne sont tracés que depuis les points
d'entrée. L'authentification est en réalité une précondition de tous les cas internes~; la relier à
chacun surchargerait la figure sans rien apprendre au lecteur. Le destinataire, lui, en est
véritablement exempt~: son suivi passe par un lien public et n'exige aucun compte.

Les deux liens \emph{«~extend~»} signalent des cas facultatifs. \emph{Optimiser l'ordre des arrêts}
étend \emph{Planifier une tournée}, parce que le dispatcheur peut toujours ordonner ses arrêts
lui-même. \emph{Encaisser un paiement} étend \emph{Livrer une commande}, parce que toutes les
commandes ne sont pas payables à la livraison.

\subsection{Diagramme de classes global}

Le diagramme de la figure~\ref{fig:class-global} présente les concepts manipulés par le système et
leurs relations. Il fixe le vocabulaire employé dans tout le reste du rapport.

\begin{figure}[H]
\centering
\includegraphics[width=\textwidth,height=0.86\textheight,keepaspectratio]{class-domain}
\caption{Diagramme de classes global}
\label{fig:class-global}
\end{figure}

Il se lit de haut en bas, du référentiel vers l'exécution. La partie haute regroupe ce qui existe
avant toute commande~: dépôts, chauffeurs, véhicules et zones. Au centre, une \textbf{Commande}
issue du progiciel se décompose en une ou plusieurs \textbf{Livraisons}~; le losange plein indique
une composition, c'est-à-dire que les livraisons n'ont pas d'existence propre en dehors de leur
commande. En parallèle, un chauffeur exécute des \textbf{Tournées}, elles-mêmes composées
d'\textbf{Arrêts}.

La classe \emph{Arrêt} est abstraite, et c'est le point le moins évident du modèle. Un arrêt de
retrait désigne un dépôt où le chauffeur charge~; un arrêt de livraison désigne une livraison à
remettre. Ce ne sont pas les mêmes objets, d'où la généralisation~: les représenter par un simple
attribut « type » aurait masqué une différence de structure réelle.

La partie basse porte les traces laissées par l'exécution~: une \textbf{PreuveDeLivraison}, un
\textbf{Encaissement} lorsque le client règle à la porte, et un \textbf{Retour} si la marchandise
revient. Les encaissements d'un chauffeur sont ensuite regroupés dans une \textbf{RemiseDeCaisse},
ce qui matérialise la séparation des responsabilités~: celui qui déclare la somme n'est pas celui
qui la valide.

Enfin, la classe \textbf{Entreprise} n'est reliée à rien, et c'est délibéré. L'appartenance d'une
donnée à une entreprise cliente n'est pas portée par une clé étrangère mais par un schéma de base de
données distinct. Le chapitre~\ref{ch:archi} détaille ce mécanisme.

\section{Gestion du projet avec Scrum}

\subsection{Product Backlog}

Le \emph{product backlog} traduit les besoins fonctionnels en éléments réalisables et priorisés. Ils
ont été formulés sous forme de \emph{user stories}, selon le gabarit « en tant que \emph{rôle}, je
veux \emph{action}, afin de \emph{bénéfice} », puis regroupés en huit épopées correspondant aux
domaines de la section précédente.

Deux échelles complètent chaque élément. La \textbf{priorité} suit la méthode MoSCoW —
\emph{Must~have} pour l'indispensable, \emph{Should~have} pour l'important non bloquant,
\emph{Could~have} pour le souhaitable. L'\textbf{estimation} est exprimée en points de complexité
selon la suite de Fibonacci~; il s'agit d'une complexité relative entre éléments, non d'une durée.

\begin{longtable}{L{1.1cm}L{8.2cm}C{1.5cm}C{1.1cm}C{1.3cm}}
\caption{Product Backlog du projet}
\label{tab:backlog} \\
\hline
\rowcolor{isimsblue!10}
\textbf{Réf.} & \textbf{User story} & \textbf{Priorité} & \textbf{Est.} & \textbf{Sprint} \\
\hline
\endfirsthead
\multicolumn{5}{c}{\tablename\ \thetable{} — \emph{suite}} \\
\hline
\rowcolor{isimsblue!10}
\textbf{Réf.} & \textbf{User story} & \textbf{Priorité} & \textbf{Est.} & \textbf{Sprint} \\
\hline
\endhead
\hline
\multicolumn{5}{r}{\emph{suite page suivante}} \\
\endfoot
\hline
\endlastfoot

\multicolumn{5}{l}{\cellcolor{isimsblue!5}\textbf{Épopée 1 — Identités, autorisations et multi-tenance}} \\
\hline
US-01 & En tant qu'utilisateur, je veux m'authentifier par un compte unique afin d'accéder aux écrans
correspondant à mon rôle. & Must & 8 & 1 \\
\hline
US-02 & En tant qu'administrateur, je veux créer des comptes et leur attribuer un rôle afin de
contrôler qui accède à quoi. & Must & 5 & 1 \\
\hline
US-03 & En tant qu'administrateur, je veux que chaque opération vérifie la permission portée par le
rôle de l'utilisateur afin que personne ne puisse dépasser ses droits. & Must & 8 & 1 \\
\hline
US-04 & En tant qu'entreprise cliente, je veux que mes données soient inaccessibles aux autres
entreprises hébergées afin de garantir ma confidentialité. & Must & 13 & 1 \\
\hline

\multicolumn{5}{l}{\cellcolor{isimsblue!5}\textbf{Épopée 2 — Intégration au progiciel de gestion}} \\
\hline
US-05 & En tant que dispatcheur, je veux importer les commandes prêtes depuis le progiciel afin de ne
pas les ressaisir. & Must & 13 & 2 \\
\hline
US-06 & En tant qu'administrateur, je veux relier les champs de mon progiciel à ceux du système afin
d'adapter l'intégration à mon paramétrage. & Must & 13 & 2 \\
\hline
US-07 & En tant qu'entreprise, je veux pouvoir brancher la solution sur n'importe quel progiciel de
gestion afin de ne dépendre d'aucun éditeur. & Must & 21 & 2 \\
\hline
US-08 & En tant que dispatcheur, je veux que le statut d'une livraison remonte automatiquement dans le
progiciel afin que la facturation suive. & Must & 8 & 2 \\
\hline
US-09 & En tant que dispatcheur, je veux qu'une synchronisation interrompue soit réessayée puis
signalée afin qu'aucune information ne se perde silencieusement. & Must & 13 & 2 \\
\hline
US-10 & En tant que dispatcheur, je veux récupérer les entrepôts du progiciel comme dépôts afin
d'éviter un double référentiel. & Should & 5 & 2 \\
\hline
US-11 & En tant qu'administrateur, je veux que le système fonctionne avec plusieurs versions du même
progiciel afin de ne pas imposer de montée de version au client. & Should & 13 & 2 \\
\hline

\multicolumn{5}{l}{\cellcolor{isimsblue!5}\textbf{Épopée 3 — Référentiel logistique}} \\
\hline
US-12 & En tant qu'administrateur, je veux gérer la liste des chauffeurs afin de refléter l'équipe
réelle. & Must & 5 & 1 \\
\hline
US-13 & En tant qu'administrateur, je veux gérer le parc de véhicules avec leur capacité afin de
pouvoir les affecter. & Should & 5 & 1 \\
\hline
US-14 & En tant qu'administrateur, je veux définir des zones géographiques afin de regrouper les
livraisons par secteur. & Should & 8 & 1 \\
\hline

\multicolumn{5}{l}{\cellcolor{isimsblue!5}\textbf{Épopée 4 — Planification et affectation}} \\
\hline
US-15 & En tant que dispatcheur, je veux construire une tournée à partir de livraisons afin
d'organiser la journée. & Must & 13 & 3 \\
\hline
US-16 & En tant que dispatcheur, je veux affecter un chauffeur et un véhicule à une tournée afin d'en
désigner le responsable. & Must & 8 & 3 \\
\hline
US-17 & En tant que dispatcheur, je veux obtenir un ordre de passage optimisé afin de réduire la
distance parcourue. & Should & 13 & 3 \\
\hline
US-18 & En tant que dispatcheur, je veux valider une tournée afin de la rendre visible au chauffeur.
& Must & 5 & 3 \\
\hline
US-19 & En tant que dispatcheur, je veux connaître l'heure d'arrivée estimée à chaque arrêt afin de
détecter les retards. & Should & 8 & 3 \\
\hline

\multicolumn{5}{l}{\cellcolor{isimsblue!5}\textbf{Épopée 5 — Exécution terrain}} \\
\hline
US-20 & En tant que chauffeur, je veux consulter ma tournée du jour sur mon téléphone afin de savoir
où aller. & Must & 8 & 4 \\
\hline
US-21 & En tant que chauffeur, je veux déclarer une livraison réussie ou échouée avec un motif afin
que le bureau soit informé. & Must & 8 & 4 \\
\hline
US-22 & En tant que chauffeur, je veux enregistrer une preuve de livraison avec photographies afin de
couvrir l'entreprise en cas de litige. & Must & 13 & 4 \\
\hline
US-23 & En tant que chauffeur, je veux continuer à travailler sans réseau et voir mes actions
transmises au retour de la couverture afin de ne jamais être bloqué. & Must & 21 & 4 \\
\hline
US-24 & En tant que chauffeur, je veux transmettre un colis à un collègue avec confirmation afin que
la responsabilité soit tracée. & Should & 13 & 4 \\
\hline

\multicolumn{5}{l}{\cellcolor{isimsblue!5}\textbf{Épopée 6 — Encaissement et remise de caisse}} \\
\hline
US-25 & En tant que chauffeur, je veux enregistrer le montant encaissé et son mode de règlement afin
d'en garder la trace. & Must & 8 & 4 \\
\hline
US-26 & En tant que chauffeur, je veux déclarer le total remis en fin de tournée afin de clore ma
journée. & Must & 5 & 4 \\
\hline
US-27 & En tant que dispatcheur, je veux compter la somme reçue et arbitrer l'écart afin que la
caisse soit vérifiable. & Must & 8 & 4 \\
\hline

\multicolumn{5}{l}{\cellcolor{isimsblue!5}\textbf{Épopée 7 — Suivi, aléas et retours}} \\
\hline
US-28 & En tant que dispatcheur, je veux voir l'avancement des livraisons sans recharger la page afin
de réagir immédiatement. & Must & 13 & 3 \\
\hline
US-29 & En tant que dispatcheur, je veux être alerté des retards et des échecs afin de traiter les
exceptions en priorité. & Must & 8 & 3 \\
\hline
US-30 & En tant que dispatcheur, je veux réaffecter une livraison bloquée à un autre chauffeur afin
de sauver la journée. & Must & 8 & 3 \\
\hline
US-31 & En tant que dispatcheur, je veux approuver ou refuser une demande de retour et déclencher
la reprise en stock afin de garder la main sur ce qui revient. & Should & 8 & 3 \\
\hline
US-32 & En tant que destinataire, je veux suivre mon colis par un lien reçu afin de connaître son
avancement sans créer de compte. & Should & 8 & 3 \\
\hline

\multicolumn{5}{l}{\cellcolor{isimsblue!5}\textbf{Épopée 8 — Pilotage et exploitation}} \\
\hline
US-33 & En tant que responsable, je veux consulter les indicateurs de performance afin de piloter
l'activité. & Should & 13 & 5 \\
\hline
US-34 & En tant que responsable, je veux exporter un rapport de tournée afin de le transmettre ou de
l'archiver. & Should & 8 & 5 \\
\hline
US-35 & En tant qu'administrateur, je veux une console indiquant l'état de l'intégration afin de
distinguer une panne du progiciel d'une panne du système. & Must & 13 & 5 \\
\hline
US-36 & En tant qu'administrateur, je veux relancer une synchronisation en échec afin de réparer sans
intervention technique. & Must & 8 & 5 \\
\hline
US-37 & En tant qu'administrateur, je veux consulter le journal des opérations sensibles afin de
pouvoir répondre à un audit. & Should & 8 & 5 \\
\hline
US-38 & En tant qu'utilisateur, je veux interroger une assistance interne sur le fonctionnement du
système afin de trouver une réponse sans quitter l'application. & Could & 13 & 5 \\
US-39 & En tant que destinataire, je veux demander le retour d'un article reçu afin de ne pas avoir
à téléphoner. & Should & 5 & 3 \
US-40 & En tant que chauffeur, je veux collecter un retour chez le destinataire afin de le ramener
au dépôt. & Should & 5 & 4 \\
\hline
\end{longtable}

Le backlog totalise quarante éléments pour 387 points de complexité. Il a été tenu dans
\textbf{GitLab Issues}, chaque épopée correspondant à une étiquette et chaque sprint à un tableau
kanban dédié.

\subsection{Planification des sprints}

Le projet s'est déroulé du 2~mars au 17~août 2026, soit vingt-quatre semaines. Il a été découpé en un
sprint d'initialisation et cinq sprints fonctionnels de trois à cinq semaines, chacun clos par une
revue et une rétrospective avant l'engagement du suivant.

L'ordre des sprints obéit à une contrainte simple~: chaque sprint ne doit dépendre que de ce qui a
déjà été livré. L'isolation par entreprise et le contrôle des droits viennent en premier, parce que
toute donnée créée ensuite doit déjà être rattachée à un client et protégée — les rajouter après coup
aurait imposé de reprendre l'ensemble du code écrit entre-temps. L'intégration au progiciel vient
juste après, parce qu'elle alimente le système en commandes~: sans elle, les sprints suivants
auraient travaillé sur des données saisies à la main, donc sur un cas qui n'existe pas en
production.

\begin{table}[H]
\centering
\caption{Planification des sprints}
\label{tab:planning}
% Colonne « Objectif » élargie et intitulés raccourcis pour qu'aucune ligne ne passe sur deux
% lignes : le tableau tient alors dans la place laissée par le texte au lieu d'être repoussé.
\small
\begin{tabular}{L{1.5cm}L{5.6cm}C{3.2cm}C{1.5cm}C{1.4cm}}
\hline
\rowcolor{isimsblue!10}
\textbf{Sprint} & \textbf{Objectif} & \textbf{Période} & \textbf{Durée} & \textbf{Points} \\
\hline
Sprint 0 & Cadrage, besoins et étude technique & 2 — 20 mars & 3 sem. & — \\
\hline
Sprint 1 & Identité, autorisation et multi-tenance & 23 mars — 17 avril & 4 sem. & 34 \\
\hline
Sprint 2 & Intégration ERP agnostique & 20 avril — 22 mai & 5 sem. & 86 \\
\hline
Sprint 3 & Planification et suivi temps réel & 25 mai — 19 juin & 4 sem. & 115 \\
\hline
Sprint 4 & Application chauffeur et hors ligne & 22 juin — 17 juillet & 4 sem. & 84 \\
\hline
Sprint 5 & Exploitation, résilience et assistance & 20 juillet — 7 août & 3 sem. & 63 \\
\hline
\multicolumn{2}{l}{Recette, rédaction et soutenance} & 10 — 17 août & 1 sem. & — \\
\hline
\end{tabular}
\end{table}

Les sprints~2 et~3 concentrent à eux seuls plus de la moitié des points. Le premier parce que
l'indépendance vis-à-vis du progiciel exige une couche d'abstraction complète avant le moindre gain
visible~; le second parce qu'il porte le cœur métier de la solution.

\section{Étude technique}

\subsection{Environnement matériel}

Le développement a été mené sur un poste unique, qui a également hébergé l'ensemble des services
conteneurisés — bases de données, courtier de messages, serveur d'identité, moteur de calcul
d'itinéraires et instances de progiciel de test. Cette dernière contrainte explique le
dimensionnement en mémoire.

\begin{table}[H]
\centering
\caption{Environnement matériel de développement}
\label{tab:materiel}
\begin{tabular}{L{4.5cm}L{10cm}}
\hline
\rowcolor{isimsblue!10}
\textbf{Élément} & \textbf{Caractéristique} \\
\hline
Processeur & Intel Core i5-12450H, 12\textsuperscript{e} génération \\
\hline
Mémoire vive & 24 Go \\
\hline
Système d'exploitation & Windows 11 Professionnel \\
\hline
Exécution des services & Docker Desktop, moteur WSL~2 \\
\hline
Essais mobiles & Appareil Android physique et émulateur \\
\hline
\end{tabular}
\end{table}

\subsection{Environnement logiciel de développement}

\begin{table}[H]
\centering
\caption{Outils de développement et de collaboration}
\label{tab:outils}
\begin{tabular}{L{4.5cm}L{10cm}}
\hline
\rowcolor{isimsblue!10}
\textbf{Outil} & \textbf{Usage dans le projet} \\
\hline
Visual Studio Code & Éditeur principal pour le code Java, TypeScript et Dart \\
\hline
Android Studio & Compilation, émulation et débogage de l'application chauffeur \\
\hline
Git et GitLab & Gestion de versions, suivi du backlog par issues et tableaux, intégration continue \\
\hline
Docker Desktop & Exécution locale de l'ensemble des services et de leurs dépendances \\
\hline
Postman & Essais manuels des interfaces de programmation \\
\hline
DBeaver & Inspection des schémas et des données PostgreSQL \\
\hline
draw.io & Conception des diagrammes UML du présent rapport \\
\hline
\LaTeX{} (XeLaTeX) & Rédaction du rapport \\
\hline
\end{tabular}
\end{table}

\subsection{Technologies retenues}

Le choix des technologies a été arrêté au terme du sprint~0, chaque brique étant retenue pour une
raison rattachable à un besoin de la section~\ref{sec:besoins}. Le chapitre~\ref{ch:archi} en détaille
la mise en œuvre.

\begin{table}[H]
\centering
\caption{Technologies retenues et motif du choix}
\label{tab:technos}
\begin{tabular}{L{2.5cm}L{3.4cm}L{8.6cm}}
\hline
\rowcolor{isimsblue!10}
\textbf{Couche} & \textbf{Technologie} & \textbf{Motif du choix} \\
\hline
Services & Spring Boot 3.3.5, Java 17 & Écosystème mature pour les architectures en services, avec
prise en charge native de la sécurité par jetons et du multi-schéma \\
\hline
Passerelle & Spring Cloud Gateway & Point d'entrée unique~: validation des jetons et application des
droits avant tout accès aux services \\
\hline
Interface web & React 19, Vite, TypeScript & Rendu réactif adapté à un tableau de dispatch qui se met
à jour en continu \\
\hline
Application mobile & Flutter & Base de code unique, et stockage local performant sur lequel repose le
mode hors ligne \\
\hline
Base de données & PostgreSQL 16 & Schéma par entreprise cliente pour l'isolation, type
\texttt{jsonb} pour les données de forme variable, absence de coût de licence \\
\hline
Messagerie & RabbitMQ 4 & Découplage des échanges avec le progiciel, relances et file de rebut
intégrées \\
\hline
Identité & Keycloak & Serveur OpenID Connect éprouvé, gestion native des organisations \\
\hline
Fichiers & MinIO & Stockage compatible S3 des photographies de preuve, déployable localement \\
\hline
Itinéraires & OSRM & Moteur de calcul auto-hébergé, sans coût par requête ni dépendance à un service
externe \\
\hline
\end{tabular}
\end{table}

\paragraph*{Sur le choix de la base de données.} PostgreSQL a été retenu pour deux raisons directement
liées aux besoins structurants. La première est l'isolation par entreprise cliente~: PostgreSQL
permet de créer un schéma par client au sein d'une même base et d'en changer dynamiquement à chaque
requête, mécanisme sur lequel repose toute la multi-tenance décrite au chapitre~\ref{ch:archi}. La
seconde est le stockage semi-structuré~: le détail des lignes de commande et les champs
supplémentaires propres à chaque client varient d'un progiciel à l'autre, ce que le type
\texttt{jsonb} absorbe sans multiplier les tables. S'y ajoute l'absence de coût de licence, cohérente avec une solution
destinée à être déployée chez plusieurs clients.

\section{Conclusion}

Ce sprint d'initialisation a produit quatre livrables~: la liste des acteurs et de leurs droits, la
spécification des besoins fonctionnels et non fonctionnels, les deux diagrammes qui fixent le
vocabulaire du système, et un backlog de quarante éléments réparti sur cinq sprints.

Son apport principal n'est pourtant pas cette production documentaire. Il tient à trois exigences que
l'analyse a fait remonter au rang de contraintes structurantes~: l'ouverture à plusieurs progiciels,
l'hébergement de plusieurs entreprises clientes sur une même installation, et la continuité de
service en l'absence de réseau. Aucune n'est une fonctionnalité visible pour l'utilisateur, et
pourtant toutes trois ont déterminé l'architecture, l'ordre des sprints et une partie des choix
technologiques.

Le chapitre suivant présente cette architecture. Les cinq chapitres qui lui succèdent déroulent
ensuite les sprints dans l'ordre du tableau~\ref{tab:planning}.
